% Mohegan-Pequot (Prince-Speck 1904) — Source Workflow
% Issue #1368
% Status: Research complete, implementation pending

\section{Mohegan-Pequot: Prince-Speck 1904}
\label{sec:1368}

\subsection{Source Description}

\subsubsection{Publication History and Citation Clarification}

The glossary was published twice in slightly different venues, which causes
persistent citation confusion. The first printing appeared in \emph{American
Anthropologist} n.s.\ 6(1): 18--45 (January--March 1904). A typographically
identical reprint appeared in the \emph{Boas Anniversary Volume} (G.\ E.\
Stechert, New York, 1906), pp.\ 18--45. Because the glossary appeared first
in \emph{American Anthropologist} in 1904, that is the correct year for
scholarly citation. The 1906 reprint is typographically identical---same page
numbering, same text---which is why later references often cite the Boas
volume without noticing the earlier \emph{AA} publication.

The exact title is ``Glossary of the Mohegan-Pequot Language'' (not ``The
Mohegan-Pequot Language'' as sometimes mis-cited). The work is a glossary
with brief ethnographic introduction, not a grammar.

\subsubsection{J. Dyneley Prince (1868--1945): The Transcriber}

John Dyneley Prince was a remarkable polymath: diplomat, politician, and
academic linguist. Born in New York City in 1868, he earned his A.B.\ (1888)
and Ph.D.\ (1892) from Columbia University and also studied in Europe. He
served as Professor of Germanic Languages at Columbia (later Dean), as a New
Jersey State Senator, and as U.S.\ Minister to Denmark (1921--1926) and
Yugoslavia (1926--1932). He was reported to know approximately 30 languages
including Estonian, Turkish, Arabic, Slavic languages, and Hungarian.

Prince was trained in the \textbf{Neogrammarian (Junggrammatiker)} tradition,
which dominated Germanic linguistics in the late 19th century. This background
has specific consequences for his transcription:

\begin{itemize}
  \item \textbf{Phonetic impressionism over phonemic analysis:}
    Neogrammarians believed in recording ``exact'' phonetic detail rather
    than phonemic contrasts. Prince's transcription therefore varies between
    \textit{a}, \textit{ah}, \textit{au}, \textit{o} for what is likely a
    single phoneme /aː/ or /ɔː/, writes vowel length inconsistently, and
    records allophonic variation as if it were phonemic.
  \item \textbf{Germanic spelling conventions:} Prince unconsciously applied
    German orthographic habits even when intending to use English spelling
    (see \S\ref{sec:1368-german}).
  \item \textbf{No IPA training:} The glossary predates the International
    Phonetic Alphabet (1888/1900). Prince had no standardized phonetic
    alphabet---he used a blended English-German ad-hoc system.
\end{itemize}

\textbf{Critical distinction:} Prince was NOT a fieldworker. He transcribed
from Speck's field notes. He never heard the language spoken himself except
possibly in brief demonstration. This means every transcription is two steps
removed from the speaker: Fidelia Fielding $\rightarrow$ Speck's ears
$\rightarrow$ Speck's field notation $\rightarrow$ Prince's desk transcription
$\rightarrow$ published form.

\subsubsection{Frank G. Speck (1881--1950): The Fieldworker}

Frank G. Speck was a Boasian anthropologist whose fieldwork methodology was
ethnographic rather than strictly linguistic. Born in New York in 1881, he
studied at Columbia University under Franz Boas (Ph.D.\ 1905) and later
taught at the University of Pennsylvania. He was a prolific ethnographer of
Eastern Algonquian peoples (Penobscot, Mohegan, Nanticoke, Tutelo, etc.).

Speck conducted fieldwork with \textbf{Fidelia Fielding (Aiji, 1827--1908)},
the last fluent speaker of Mohegan-Pequot. Fielding was born in 1827 on the
Mohegan reservation in Uncasville, Connecticut, a direct descendant of the
Mohegan chief Uncas. She was the keeper of Mohegan traditions, maintained a
diary in Mohegan using her own spelling system, was a monolingual Mohegan
speaker as a child (learning English later in life), and was known for her
fierce preservation of the language---she refused to speak English with other
Mohegans.

Fielding was literate in English and developed her own orthography for
recording Mohegan, primarily in her personal diaries. Her spelling system
was based on English conventions as she understood them, with considerable
idiosyncrasy.

\subsubsection{The Speck--Prince Handoff}

The exact chain of custody for the data is:

\begin{enumerate}
  \item Fidelia Fielding produces utterances in Mohegan
  \item Frank Speck hears and records them in field notation
  \item Speck provides his field notes to J.\ Dyneley Prince
  \item Prince converts Speck's notation into the published orthography
  \item Prince writes the introductory notes and publishes under both names
\end{enumerate}

This two-stage recording (fieldworker $\rightarrow$ transcriber) introduces
compounding transcription errors: Speck's hearing of Fielding's speech
(accent, speed, dialect), Speck's field notation choices (ad-hoc English-based
spelling), Prince's interpretation of Speck's notation (he never heard the
original), Prince's application of German orthographic conventions, and
typesetting errors in the 1904 publication.

\subsubsection{Fidelia Fielding's Orthography and Its Influence}

Prince and Speck explicitly state that they sometimes deferred to Fielding's
spelling for ``sentimental reasons''---i.e., they kept her English-based
spelling even when a more phonetic transcription would have been technically
better. From analyses of her diaries, Fielding's orthography shows:

\begin{itemize}
  \item English vowel letters for approximate sounds (\textit{a, e, i, o, u}
    as in English orthography) with no consistent length marking
  \item \textit{gh} for /k/ or /x/, following English convention, which may
    reflect actual velar fricative
  \item \textit{w} for /w/ and as glide, consistent with English
  \item Doubled consonants for occasional geminate marking, which may reflect
    actual length or stress
  \item No systematic nasalization marking---nasal vowels were likely present
    but unwritten
\end{itemize}

The Prince-Speck preface indicates they intentionally preserved some of
Fielding's spellings because she was the last speaker and her written form
had cultural value. This means some words in the glossary may be Fielding's
\textbf{written} spelling rather than Prince's \textbf{hearing-based}
transcription, and the distinction between ``written by Fielding'' and
``recorded by Speck/Prince'' is not marked in the glossary.

\subsubsection{Orthography System}

The glossary presents a three-way orthographic competition:

\begin{table}[h]
\centering
\begin{tabular}{llll}
\toprule
\textbf{Source} & \textbf{Convention} & \textbf{Example} & \textbf{Interpretation} \\
\midrule
English colonial & \textit{ch} for /tʃ/, \textit{sh} for /ʃ/ & ``cheese'' & Familiar, predictable \\
English vernacular (Fielding) & Ad-hoc English-based & Variable spelling & Unpredictable \\
German linguistic (Prince) & \textit{tsh} for /tʃ/, \textit{z} for /ts/, Dehnungs-\textit{h} & Specific forms & Detectable pattern \\
\bottomrule
\end{tabular}
\caption{Orthographic confusion index}
\label{tab:1368-orthographic}
\end{table}

\subsection{Phonemic Reconstruction}
\label{sec:1368-phonemic}

\subsubsection{Consonant Inventory (Reconstructed)}

\begin{table}[h]
\centering
\begin{tabular}{lccccc}
\toprule
 & Bilabial & Alveolar & Palatal & Velar & Glottal \\
\midrule
Stops & p & t & & k & \\
Voiced stops & b? & d? & & g? & \\
Affricates & & & tʃ & & \\
Fricatives & & s & ʃ & (x?) & h \\
Nasals & m & n & & & \\
Approximants & w & & j & & \\
\bottomrule
\end{tabular}
\caption{Reconstructed consonant inventory for Mohegan-Pequot}
\label{tab:1368-consonants}
\end{table}

Notes: \textit{b, d, g} are likely English loan phenomena or allophones of
/p, t, k/. The velar fricative /x/ may exist but is not clearly distinguished
from /h/ or /k/ in the transcription. Preaspiration /ʰp, ʰt, ʰk/ is probably
a suprasegmental feature of the syllable coda, not a phoneme sequence.

\subsubsection{Vowel Inventory (Reconstructed)}

\begin{table}[h]
\centering
\begin{tabular}{lccc}
\toprule
 & Front & Central & Back \\
\midrule
High & i /iː/ & & u /uː/ \\
Mid & e /ɛ/ & ə /ə/ & o /o/ \\
Low & & a /a/ & ɔ /ɔː/ \\
\bottomrule
\end{tabular}
\caption{Reconstructed vowel inventory for Mohegan-Pequot}
\label{tab:1368-vowels}
\end{table}

Diphthongs: /aj/ (spelled \textit{ai}), /aw/ (spelled \textit{au}), /oj/
(if present).

\subsubsection{Preaspiration Phonology}

Preaspiration in SNEA languages is a debated feature. The Prince-Speck
evidence suggests that preaspiration occurred before fortis stops in coda
position: \textit{ahk}, \textit{ahp}, \textit{aht}. It may have been a
syllable-level prosodic feature (fortis/lenis distinction) rather than a
segment. Prince-Speck writes it as \textit{hC} sequences, but these may
represent single preaspirated segments /ʰC/. The fact that it appears in
Prince-Speck and not in earlier colonial sources may mean: (a) it developed
later in Mohegan-Pequot, (b) earlier recorders missed it, or (c) Prince's
Germanic ear was better tuned to it.

\subsubsection{Nasalization}

Mohegan-Pequot had nasal vowels (inherited from Proto-Algonquian *ē, *ā,
etc.). Prince-Speck does not systematically mark nasalization, may
occasionally indicate it through vowel quality choices (e.g., \textit{o} for
nasalized /ã/), and may miss it entirely, mapping nasal vowels to oral
vowels. Fielding's Mohegan has lost nasal vowels, so the modern source is no
help for this feature.

\subsection{German Orthographic Influence}
\label{sec:1368-german}

Prince's Germanistic training is visible throughout the transcription.

\subsubsection{Affricate Representations}

\begin{table}[h]
\centering
\begin{tabular}{llll}
\toprule
\textbf{Prince-Speck} & \textbf{German convention} & \textbf{English equivalent} & \textbf{Phoneme} \\
\midrule
\textit{tsh} & German \textit{tsch} (as in \textit{deutsch}) & Rare, used in ``batsha'' & /tʃ/ \\
\textit{z} & German \textit{z} = /ts/ & English \textit{z} = /z/ & ambiguous \\
\bottomrule
\end{tabular}
\caption{Affricate representations}
\label{tab:1368-affricates}
\end{table}

The use of \textit{tsh} for the /tʃ/ affricate is a distinctly Germanic
choice. English colonial recorders typically used \textit{ch} or \textit{tch}.
The \textit{ts} prefix represents a German-trained linguist's desire to show
the affricate's stop onset.

\subsubsection{Sibilant Representations}

\begin{itemize}
  \item \textit{sch} (if present): German \textit{sch} = /ʃ/ (English
    \textit{sh} would be more typical of American recorders)
  \item \textit{s} before vowels: German \textit{s} = /z/ in some positions
    (English reads this as /s/)
\end{itemize}

\subsubsection{Vowel Length Marking}

German linguistics of the period marked vowel length with vowel doubling
(\textit{ee}, \textit{oo})---shared with English, so the source is
ambiguous---and \textit{h} after vowel (\textit{ah}, \textit{eh})---German
uses \textit{Dehnungs-h} (lengthening h), which is NOT the same as English
``ah'' for /aː/. Prince's \textit{ah} spelling could derive from either
English ``ah'' (indicating quality as much as length) or German
\textit{Dehnungs-h} (indicating length primarily). The ambiguity cannot be
resolved without acoustic evidence.

\subsubsection{Consonant Fortition}

German-influenced recorders sometimes write initial /b, d, g/ as \textit{b,
d, g} but intend fortis articulation, apply final devoicing even when
hearing voiced finals---writing \textit{d} but hearing /t/, and use
unaspirated \textit{p, t, k} where English recorders would write \textit{b,
d, g}.

\subsubsection{Umlaut-Like Notation}

German linguists sometimes used \textit{ü}, \textit{ö}, \textit{ä} when
hearing sounds not in English or German. Whether Prince does this in the
glossary requires a full search, but it is a known pattern in
German-trained linguist transcriptions.

\subsection{Preaspiration Evidence}

\subsubsection{Explicit /h/ in Prince-Speck}

One of Prince-Speck's most valuable features for SNEA historical phonology
is its treatment of /h/:

\begin{itemize}
  \item \textbf{Word-initial h:} \textit{ahupanun} (come here)---English
    recorders often drop word-initial h
  \item \textbf{Intervocalic h:} \textit{bahduntah} (rising)---often missed
    by colonial recorders
  \item \textbf{Cluster hC:} \textit{bahkeder} (bitter?)---PREASPIRATION,
    critical for SNEA reconstruction
\end{itemize}

\subsubsection{Preaspiration (/hk/, /hp/, /ht/)}

The spelling \textit{hk} in \textit{bahkeder} is significant because it
represents a phonetic feature that English-colonial recorders systematically
miss. Roger Williams (Narragansett) writes plain \textit{k} for both /k/ and
/hk/. John Eliot (Wampanoag) uses doubled vowels to indicate syllable weight
but does not write preaspiration. Ezra Stiles (Mohegan-Pequot) sometimes
writes \textit{hk} but inconsistently. \textbf{Prince writes \textit{hk},
\textit{hp}, \textit{ht} clusters explicitly}---this is unique among the
pre-modern SNEA sources.

\subsubsection{Why Prince-Speck Catches Preaspiration}

Prince's Germanic training gave him an ear for aspiration contrasts. German
distinguishes aspirated /pʰ, tʰ, kʰ/ in certain positions from unaspirated
/p, t, k/ in clusters, and the \textit{Dehnungs-h} convention made him
conscious of h-like elements. Additionally, Mohegan-Pequot may have preserved
preaspiration longer than other SNEA languages, or Speck's field notes may
have captured it from Fidelia Fielding's careful speech.

\subsubsection{Comparison with Other Sources}

\begin{table}[h]
\centering
\begin{tabular}{llll}
\toprule
\textbf{Source} & \textbf{Preaspiration /hk/} & \textbf{Intervocalic /h/} & \textbf{Word-initial /h/} \\
\midrule
Williams 1643 & Not written & Sometimes & Rarely \\
Eliot 1660s & Not written & Rarely & Rarely \\
Stiles 1760s--90s & Inconsistent & Inconsistent & Inconsistent \\
\textbf{Prince-Speck 1904} & \textbf{Written} & \textbf{Written} & \textbf{Written} \\
Fielding 2013 & Not phonemic & Not present & Not present \\
\bottomrule
\end{tabular}
\caption{Preaspiration and /h/ recording across SNEA sources}
\label{tab:1368-preaspiration}
\end{table}

This makes Prince-Speck a \textbf{key witness for preaspiration} in SNEA
historical phonology, on par with Mayhew and Bourne (Wampanoag) for the same
feature.

\subsection{English Loanword Typology}

The Prince-Speck glossary contains a very high proportion of English
loanwords, reflecting the contact situation of early 20th-century Mohegan.

\subsubsection{Loanword Categories}

\textbf{Category 1: Cultural/Technological (new concepts)}

\begin{itemize}
  \item \textit{beed} (bed), \textit{beeddunk} (bedstead), \textit{beesh}
    (peas), \textit{beetkuz} (lady's dress), \textit{manshetshe} (matches),
    \textit{tchiken} (chicken), \textit{breeches} (breeches), \textit{cock}
    (rooster), \textit{nute} (nut), \textit{rug} (rug), \textit{soyler}
    (soldier), \textit{spuner} (spoon), \textit{shtoo} (stew), \textit{tappe}
    (tap/faucet), \textit{tool} (tool), \textit{wamp} (wampum, earlier
    borrowing)
\end{itemize}

\textbf{Category 2: Extended/Replaced Domains}

\begin{itemize}
  \item \textit{lif} (leaf---English ``leaf'' replacing native word),
    \textit{appece} (apple---meaning extended to any fruit), \textit{cheese}
    (cheese), \textit{beksees} (pig---English ``pigs'' with Mohegan suffix?),
    \textit{square} (square)
\end{itemize}

\textbf{Category 3: Calendric/Temporal}

\begin{itemize}
  \item \textit{beitar} (Friday---possibly from English ``Friday'' via Dutch
    ``Vrijdag''?)
\end{itemize}

\textbf{Category 4: Phonological Nativization}

\begin{itemize}
  \item \textit{manshetshe} (matches): /tʃ/ $\rightarrow$ \textit{tsh};
    final vowel epenthesis
  \item \textit{soyler} (soldier): r $\rightarrow$ /l/ (no native /r/);
    final /ər/ $\rightarrow$ /er/
  \item \textit{spuner} (spoon): no initial cluster simplification;
    final /n/ $\rightarrow$ /nər/
  \item \textit{appece} (apple): /l/ $\rightarrow$ ?; final schwa
    $\rightarrow$ \textit{e}
\end{itemize}

\subsubsection{Identifying Loanwords}

Loanwords in the glossary can be identified by phonotactic cues (presence of
/r/, /l/, initial /b,d,g/ without alternation, consonant clusters not typical
of Algonquian), semantic cues (European cultural artifacts, introduced
concepts, calendric terms), and orthographic cues (English spelling patterns
that diverge from Mohegan-Pequot patterns).

\subsubsection{Handling in Conversion Pipeline}

Loanwords should be kept with identity mapping (loanword $\rightarrow$ same
word) rather than pushed through Mohegan-Pequot phoneme conversion. The
detection FSM should flag suspected loanwords at runtime based on phonotactic
patterns and semantic context.

\subsection{Comparison with Other SNEA Recording Traditions}

\subsubsection{Williams 1643 (Narragansett)}

\begin{itemize}
  \item /h/ recording: Williams poor, Prince-Speck good---Prince better for
    preaspiration
  \item Vowel system: Williams English short/long, Prince-Speck more
    detailed---both need normalization
  \item Loanword handling: Williams marked as loans, Prince-Speck
    unmarked---Prince needs loanword detection
  \item Transcription consistency: Williams moderate, Prince-Speck low
    (two-stage)---Prince less consistent
  \item Grammar coverage: Williams extensive, Prince-Speck minimal---Prince
    is glossary only
\end{itemize}

\subsubsection{Eliot 1660s (Wampanoag Mass Bible)}

\begin{itemize}
  \item Orthography: Eliot semi-standardized, Prince-Speck ad-hoc
  \item Corpus size: Eliot massive (~1000pp), Prince-Speck small (446 words)
  \item Phonemic value: Eliot low (English-colonial), Prince-Speck
    potentially high (German training)
  \item Preaspiration: Eliot not written, Prince-Speck written
\end{itemize}

\subsubsection{Stiles 1760s--90s (Mohegan-Pequot)}

Ezra Stiles (1727--1795), president of Yale, collected Mohegan-Pequot
vocabulary from several speakers (including Occom family members) over
decades.

\begin{itemize}
  \item Time depth: Stiles ~140 years earlier, Prince-Speck later, closer to
    extinction
  \item Number of speakers: Stiles multiple (including Occom), Prince-Speck
    single (Fidelia Fielding)
  \item Transcription: Stiles ad-hoc English, Prince-Speck Germanic-influenced
  \item /h/ and preaspiration: Stiles inconsistent, Prince-Speck explicit
  \item Dialect: both Mohegan-Pequot proper
  \item Loanword rate: Stiles lower, Prince-Speck higher (more contact)
\end{itemize}

Stiles and Prince-Speck together provide diachronic depth---comparing the
same language approximately 140 years apart.

\subsubsection{Comparison with Stephanie Fielding (2013) Modern Mohegan}

Stephanie Fielding (Yôpôwi Yoht, ``Morning Fire''), a Mohegan linguist
educated at MIT, published the \emph{Modern Mohegan Dictionary} in 2013
(revised 2015). This is the only modern comprehensive dictionary of Mohegan.

\begin{table}[h]
\centering
\begin{tabular}{lll}
\toprule
\textbf{Feature} & \textbf{Fielding (2013)} & \textbf{Prince-Speck (1904)} \\
\midrule
Base & Phonemic (designed) & Phonetic (impressionistic) \\
Vowel length & Explicit & Inconsistent \\
/h/ & Not phonemic & Orthographic \\
Preaspiration & Not marked & Marked \\
Nasalization & Not marked & Not marked \\
Loanwords & Identified & Unidentified \\
Consonant system & Simplified (modern) & More complex (archaic) \\
\bottomrule
\end{tabular}
\caption{Orthography comparison: Fielding vs.\ Prince-Speck}
\label{tab:1368-fielding-comparison}
\end{table}

The pipeline of Prince-Speck $\rightarrow$ Fielding $\rightarrow$ Proto-SNEA
is powerful for reconstruction. Fielding gives the modern phonemic form (what
the language sounds like now); Prince-Speck gives the older phonetic form
(what the language sounded like ~110 years ago); differences between them
indicate sound changes in progress. However, Fielding's dictionary is
reconstructive---many entries are not from fluent speakers but from archival
sources (including Prince-Speck). This introduces circularity: Fielding may
have reconstructed a word FROM Prince-Speck, so using Fielding to ``verify''
Prince-Speck is invalid for that word.

\subsection{Problems Encountered}

\begin{itemize}
  \item \textbf{German orthographic influence:} Prince's Neogrammarian
    training imposed German conventions: \textit{tsh} for /tʃ/, \textit{z}
    for /ts/, \textit{Dehnungs-h} for vowel length, final devoicing, and
    possible umlaut-like notation for non-English vowels.
  \item \textbf{English loanword density:} Approximately 20--30\% of entries
    are English loanwords (bed, matches, chicken, spoon, soldier). These must
    be detected and excluded from phoneme conversion.
  \item \textbf{Fielding dictionary circularity:} Many Fielding (2013) entries
    are reconstructed FROM Prince-Speck---using Fielding to verify
    Prince-Speck is circular for those entries. A cross-reference is needed
    to identify which Fielding entries derive from Prince-Speck
    (self-source) vs.\ from other sources (independent verification).
  \item \textbf{Record count correction:} 446 entries, not 465 as previously
    stated in some references. The actual published glossary contains 446
    numbered entries (some are phrases rather than single words).
  \item \textbf{Two-stage recording chain:} Speck collected, Prince
    transcribed---Prince never heard the language spoken. This introduces
    compounding transcription errors at each stage.
  \item \textbf{``Sentimental reasons'' spellings:} Prince and Speck
    intentionally preserved some of Fielding's English-based spellings for
    cultural value, making it impossible to distinguish Fielding's written
    spelling from Prince's hearing-based transcription.
  \item \textbf{Inconsistent vowel length marking:} Prince writes vowel
    length sometimes with \textit{h}, sometimes with doubling, sometimes not
    at all---the source of the length marking (English ``ah'' vs.\ German
    \textit{Dehnungs-h}) is ambiguous.
  \item \textbf{Multi-word entries:} The glossary includes multi-word phrases,
    which complicates conversion because phrases span different syntactic
    contexts.
\end{itemize}

\subsection{Research Approach}

The conversion pipeline follows this sequence:

\begin{enumerate}
  \item \textbf{Loanword detection FSM:} Automated phonotactic scan of all
    446 entries to identify and catalog English loanwords. Loanwords are
    flagged for identity mapping (excluded from phoneme conversion).
  \item \textbf{German orthography normalizer:} Identify all words where
    German conventions (\textit{tsh}, \textit{z}, \textit{Dehnungs-h})
    appear and apply the correct phoneme mapping.
  \item \textbf{Fielding cross-reference with circularity flag:} Map each
    Prince-Speck entry against Fielding (2013) with source tracing---identify
    which Fielding entries derive from Prince-Speck (circular) vs.\
    independent sources.
  \item \textbf{Preaspiration capture:} Extract all \textit{hC} clusters
    from the glossary and classify by position (initial, medial, final).
    Prince-Speck is one of the few sources that writes preaspiration
    (/hC/ clusters) explicitly---this must be preserved.
  \item \textbf{Stiles comparison:} Cross-reference Prince-Speck entries
    against Stiles's Mohegan-Pequot vocabulary for diachronic comparison.
  \item \textbf{Phoneme-to-orthography mapping table:} Full conversion table
    from Prince-Speck spelling to phonemes, covering all 446 entries.
\end{enumerate}

\subsection{Conversion Workflow}
\label{sec:1368-workflow}

\texttt{[Pending implementation --- see Issue \#1368]}

The conversion workflow will implement the research approach described above
as a multi-stage pipeline. Each stage will be implemented as a separate
module with its own test suite.

\subsection{Linguist Feedback}
\label{sec:1368-feedback}

\texttt{[Template --- content pending review]}

The following aspects require linguist review:

\begin{itemize}
  \item German orthography pattern extraction: confirmation of detected
    German conventions
  \item Preaspiration audit: verification of /hC/ cluster classification
  \item Loanword catalog: review of flagged loanwords for accuracy
  \item Fielding cross-reference: validation of circularity detection
  \item Phoneme mapping: review of the conversion table for accuracy
\end{itemize}

\subsection{Related Work: Colonial Recorder Perception Problem}

The core challenge of this research card---English-speaking recorders
misperceiving and mis-transcribing Algonquian phonemes---is a known
phenomenon documented across Eastern North America. Two researchers' work
is directly relevant.

\textbf{Dr.\ Keith Cunningham} (Linguist, Nanticoke Indian Tribe). His
doctoral dissertation \emph{A Phonological Analysis of Nanticoke With
Practical Applications for Language Revitalization} (Georgetown University)
analyzes three 18th-century Nanticoke word lists (Heckewelder 1785) using
the same methodology: historical phonology from colonial recorders, with the
same challenges of multilingual informants and orthographic complexity. His
2024 Algonquian Conference presentation \emph{A Revised Phonological Analysis
of the Heckewelder Vocabulary of Nanticoke} directly parallels the
Prince-Speck analysis---both involve a two-stage recorder chain (fieldworker
$\rightarrow$ desk transcriber) that introduces compounding transcription
errors, and both require identifying the transcriber's native-language
orthographic conventions (German for Prince, English for Heckewelder).
Cunningham's work demonstrates that the colonial recorder perception problem
extends beyond SNEA into the broader Eastern Algonquian family.

\textbf{Dr.\ Craig Kopris} (Independent scholar). His paper \emph{Les sons
wyandots perçus par des oreilles étrangères} (Wyandot sounds perceived by
foreign ears, \emph{Recherches amérindiennes au Québec}, 2014) is the exact
framing of the Prince-Speck problem: how European recorders systematically
misperceived and mis-transcribed indigenous phonemes. His \emph{Wyandot
Phonology: Recovering the Sound System of an Extinct Language} applies the
same recovery methodology to an unrelated language family (Iroquoian),
demonstrating that the foreign-ear perception problem is universal across
colonial documentation of North American languages. His Dictionary of
Wyandot project (NEH award FN-255576-17) shows the long-term application of
these methods to language revitalization.

The Prince-Speck glossary (446 entries) is unique among the SNEA sources in
having a documented two-stage fieldworker-transcriber chain (Speck collected,
Prince transcribed) with known linguistic biases (Prince's Neogrammarian
training). This makes it the most directly comparable case to the
recorder-chain problems that both Cunningham and Kopris address in their
respective language families.

\subsection{References}

\begin{itemize}
  \item Prince, J.\ Dyneley \& Speck, Frank G.\ (1904). ``Glossary of the
    Mohegan-Pequot Language.'' \emph{American Anthropologist} n.s.\ 6(1):
    18--45.
  \item Prince, J.\ Dyneley \& Speck, Frank G.\ (1906). ``Glossary of the
    Mohegan-Pequot Language.'' In \emph{Boas Anniversary Volume}, pp.\
    18--45. New York: G.\ E.\ Stechert. [Reprint, typographically identical]
  \item Prince, J.\ Dyneley. (1907). ``A Singular Mohegan-Pequot Form.''
    \emph{American Anthropologist} n.s.\ 9: 496--497.
  \item Fielding, Stephanie. (2013/2015). \emph{Modern Mohegan Dictionary.}
    Mohegan Tribe.
  \item Cowan, William. (1973). ``Pequot from Stiles to Speck.''
    \emph{International Journal of American Linguistics} 39(3): 164--172.
  \item Speck, Frank G.\ (1928). ``Native Tribes and Dialects of Connecticut:
    A Mohegan-Pequot Diary.'' \emph{Annual Report of the Bureau of American
    Ethnology} 43: 199--287.
  \item Goddard, Ives. (1978). ``Eastern Algonquian Languages.'' In
    \emph{Handbook of North American Indians}, Vol.\ 15: Northeast, pp.\
    70--77. Washington: Smithsonian Institution.
  \item Costa, David J.\ (2007). ``The Dialectology of Southern New England
    Algonquian.'' In \emph{Papers of the 38th Algonquian Conference}, pp.\
    81--127.
\end{itemize}
